\documentclass[11pt]{article}
\usepackage[margin=1in]{geometry}
\usepackage{graphicx}
\usepackage{amsmath}
\usepackage{array}
\usepackage{hyperref}

\title{DRX--Carbon Morphology Modeling and Compaction Simulation}
\author{Project Summary}
\date{\today}

\begin{document}
\maketitle

\section{Overview}
This project builds mesoscale particle packs for a disordered rocksalt (DRX) cathode mixed with conductive carbon and prepares discrete element method (DEM) simulations to study carbon morphology effects during compaction. All geometry is generated in micrometer-scale units to remain compatible with LIGGGHTS/LAMMPS granular models. Three carbon morphologies---Super~P aggregates, carbon nanotube (CNT) chains, and graphene flakes---complement a homogeneous $0.1\,\mu$m carbon baseline.

\section{Mesh Generation Pipelines}
Each morphology lives in its own folder under the project root. Table~\ref{tab:morph} summarizes the configurations.

\begin{table}[h]
    \centering
    \begin{tabular}{|>{\raggedright}p{3cm}|p{4cm}|p{5cm}|}
        \hline
        \textbf{Folder} & \textbf{Carbon Morphology} & \textbf{Notes} \\
        \hline
        \texttt{mesh\_superp} & Super~P clusters (0.5~$\mu$m aggregates of $0.1~\mu$m spheres) & Clusters precomputed with fixed offsets; no rotation during placement. \\
        \texttt{mesh\_cnt} & CNT rods (1~$\mu$m length, $0.1~\mu$m diameter spheres) & Chains allow random 3D orientation for isotropic connectivity. \\
        \texttt{mesh\_graphene} & Graphene flakes (1~$\mu$m~$\times$~1~$\mu$m, 0.1~$\mu$m thick) & Rectangular grids of spheres that can rotate to introduce misalignment. \\
        \texttt{mesh\_homo} & Homogeneous $0.1~\mu$m carbon spheres & Baseline without clustering; copied from the original code. \\
        \hline
    \end{tabular}
    \caption{Carbon morphology workspaces. All share the same DRX distribution and ghost-particle handling.}
    \label{tab:morph}
\end{table}

\subsection{Input Parameters}
\begin{itemize}
    \item \textbf{DRX distribution}: Log-normal, mean diameter $1~\mu$m, min--max 0.5--2~$\mu$m, density $4.0\,\text{pg}/\mu\text{m}^3$.
    \item \textbf{Carbon building block}: Sphere of diameter $0.1~\mu$m, density $2.2\,\text{pg}/\mu\text{m}^3$.
    \item \textbf{Domain}: $15\,\mu$m $\times$ $15\,\mu$m $\times$ $30\,\mu$m with periodic $x,y$; mass ratio sweep varies carbon loading from 6 to 1.
    \item \textbf{Ghost layer}: Type-3 particles sit above the pack and receive the applied compaction load.
\end{itemize}

\subsection{Workflow}
For each folder:
\begin{enumerate}
    \item Run \texttt{mdl\_run.m} in MATLAB. Eleven mass-ratio cases are generated into \texttt{massratio/mr*/mdl.data}. The script reports total DRX/carbon particle counts and cluster-placement statistics.
    \item Each \texttt{mdl.data} contains DRX spheres (type 1), carbon spheres (type 2), and ghost particles (type 3) in micrometer units with densities set for ``units micro''.
\end{enumerate}

\section{Simulation Framework}
The \texttt{simulation} folder hosts a single LIGGGHTS input file (\texttt{lmp\_mr.in}) plus helper scripts. Key features:
\begin{itemize}
    \item Parameterized variables: \texttt{datafile}, \texttt{outputdir}, \texttt{label}, and \texttt{press\_tgt}. These are passed via \texttt{-var} flags to avoid editing the input deck per run.
    \item Output management: logs and dumps land in \texttt{\$\{outputdir\}/log\_\$\{label\}.lammps} and \texttt{rst\_\$\{label\}\_*.lmp}. Pressure is monitored via ghost-force reduction.
    \item Staged compaction: initial fast press (loop~1), hold (loop~2), slow ramp (loop~3), and relaxation (loop~4), all keyed to \texttt{press\_tgt}.
\end{itemize}

\subsection{Batch Scripts}
\begin{itemize}
    \item \texttt{myjob.sh}: Loads LAMMPS via Spack and executes \texttt{lmp\_mr.in}. Required environment variables are \texttt{DATAFILE}, \texttt{OUTDIR}, \texttt{LABEL}; optional \texttt{PRESS\_TGT} and \texttt{LMP\_CMD} override defaults.
    \item \texttt{run\_all.sh}: Iterates over \texttt{mesh\_*}/\texttt{massratio/mr*/mdl.data}, launching one LIGGGHTS job per mesh. Results are stored under \texttt{simulation/results/<morph>/<mrXX>}. Run with \texttt{LMP\_CMD=lmp ./run\_all.sh}.
\end{itemize}

\section{How to Reproduce}
\subsection{Generate Meshes}
\begin{enumerate}
    \item Open MATLAB or \texttt{matlab -batch}. For Super~P, run \\ 
    \texttt{cd('.../mesh\_superp'); mdl\_run}
    \item Repeat for \texttt{mesh\_cnt}, \texttt{mesh\_graphene}, and optionally \texttt{mesh\_homo}.
    \item Verify each \texttt{massratio/mr*/mdl.data} exists and inspect reported cluster stats.
\end{enumerate}

\subsection{Run Compaction}
\begin{enumerate}
    \item For a single case, set environment variables and submit: \\ 
    \texttt{DATAFILE=/path/to/mesh\_superp/massratio/mr1/mdl.data}\\
    \texttt{OUTDIR=simulation/results/superp/mr1}\\
    \texttt{LABEL=superp	extunderscore mr1}\\
    \texttt{PRESS	extunderscore TGT=2}\\
    \texttt{sbatch --export=DATAFILE,OUTDIR,LABEL,PRESS	extunderscore TGT simulation/myjob.sh}
    \item To sweep every mesh, run \texttt{simulation/run\_all.sh}. Override \texttt{PRESS\_TGT} globally if needed.
\end{enumerate}

\section{Validation}
\begin{itemize}
    \item MATLAB \texttt{mdl\_run} completed for all morphologies and mass ratios. Progress logs show zero failed placements for Super~P, CNT, and graphene clusters.
    \item LIGGGHTS input has not been executed in this environment; please run \texttt{run\_all.sh} or submit jobs via \texttt{myjob.sh} on your target cluster.
\end{itemize}

\section{Future Extensions}
\begin{itemize}
    \item Introduce stochastic orientation weighting (e.g., CNT alignment) or graded cluster sizes.
    \item Couple the DEM outputs to electronic/ionic transport solvers for structure--property linkage.
    \item Add post-processing scripts (Python/Matlab) to compute percolation metrics from the generated meshes and DEM states.
\end{itemize}

\end{document}
